\chapter*{阅读说明}
\addcontentsline{toc}{chapter}{阅读说明}
本书将概念、图解、公式与算例放在一起，帮助初学者建立机器学习知识之间的联系。
先理解一个方法要解决的问题，再跟着图中的数据或梯度走一遍，最后回查公式与适用边界。
不必先把数学附录背完，也不必一次读完全部模型。

\begin{keyidea}[当前版本]
这是第一版定稿候选，正文采用独立矢量图与表格。
图中的数值例子配有复核脚本；本版已完成逐章内容核对与读者路线校读。
阅读时请结合公式附近的变量定义、数据划分与适用条件。
\end{keyidea}

\section*{三条阅读路线}
\begin{center}
\begin{tabularx}{\textwidth}{P{30mm}Y}
\toprule
你的当前目标 & 建议路线\\
\midrule
第一次学习建模 & 第\ref{ch:learning}章建立工作流；第\ref{ch:data}--\ref{ch:tools}章准备数据、评价与建模工具；再选第\ref{ch:linear}--\ref{ch:svm}章的经典模型，或第\ref{ch:unsupervised}章的无监督方法。\\
学习神经网络 & 先读第\ref{ch:evaluation}章的数据划分与第\ref{ch:tools}章的损失；第\ref{ch:neural}--\ref{ch:numerics}章串起前向、反向、更新与稳定训练；再按任务读CNN、文本、序列或迁移。\\
遇到具体问题回查 & 用附录\ref{ch:reference}的主题页码与方法比较找入口，再回到相应图解；形状、概率与梯度见附录\ref{ch:math}。\\
\bottomrule
\end{tabularx}
\end{center}

\section*{数学不熟悉时怎样回查}
向量、矩阵与乘法形状见附录\ref{sec:math-matrix}节；导数与梯度见附录\ref{sec:math-gradient}节；
条件概率、期望与方差见附录\ref{sec:math-probability}--\ref{sec:math-random}节。
先看对应的小算例，再回到当前主题。进入注意力结构前，若不熟悉层归一化，可查第\ref{sec:cnn-normalization}节。

\section*{每个主题怎样阅读}
先读开头的任务与定义；看图中轴、箭头、颜色和形状；对照“读图”复算小例子。
随后检查提示框中的前提和常见误解。比较方法时，保持数据划分、评价口径与计算预算可比。
图中有些曲线是明确标注的趋势示意，有些来自给出的实际递推或小数据；两者不要混同。

本书偏向可解释的小算例与机制，不提供“某算法总是最好”的选型口诀。
不同文献可能采用不同损失缩放、矩阵方向或更新符号；先对照本节定义，再换算。
“更低/更高”均需结合指标方向、样本与使用条件理解。

配套图源、算例复核脚本与构建说明见
\href{https://github.com/zhang-haifei/ML-Cheat-Sheet/blob/work/review/handbook-source.zip}{可编译源码包}；
解压后按\texttt{README.md}操作。
